% Options for packages loaded elsewhere
\PassOptionsToPackage{unicode}{hyperref}
\PassOptionsToPackage{hyphens}{url}
\documentclass[
  conference]{IEEEtran}
\usepackage{xcolor}
\usepackage{amsmath,amssymb}
\setcounter{secnumdepth}{-\maxdimen} % remove section numbering
\usepackage{iftex}
\ifPDFTeX
  \usepackage[T1]{fontenc}
  \usepackage[utf8]{inputenc}
  \usepackage{textcomp} % provide euro and other symbols
\else % if luatex or xetex
  \usepackage{unicode-math} % this also loads fontspec
  \defaultfontfeatures{Scale=MatchLowercase}
  \defaultfontfeatures[\rmfamily]{Ligatures=TeX,Scale=1}
\fi
\usepackage{lmodern}
\ifPDFTeX\else
  % xetex/luatex font selection
\fi
% Use upquote if available, for straight quotes in verbatim environments
\IfFileExists{upquote.sty}{\usepackage{upquote}}{}
\IfFileExists{microtype.sty}{% use microtype if available
  \usepackage[]{microtype}
  \UseMicrotypeSet[protrusion]{basicmath} % disable protrusion for tt fonts
}{}
\makeatletter
\@ifundefined{KOMAClassName}{% if non-KOMA class
  \IfFileExists{parskip.sty}{%
    \usepackage{parskip}
  }{% else
    \setlength{\parindent}{0pt}
    \setlength{\parskip}{6pt plus 2pt minus 1pt}}
}{% if KOMA class
  \KOMAoptions{parskip=half}}
\makeatother
\usepackage{longtable,booktabs,array}
\newcounter{none} % for unnumbered tables
\usepackage{calc} % for calculating minipage widths
% Correct order of tables after \paragraph or \subparagraph
\usepackage{etoolbox}
\makeatletter
\patchcmd\longtable{\par}{\if@noskipsec\mbox{}\fi\par}{}{}
\makeatother
% Allow footnotes in longtable head/foot
\IfFileExists{footnotehyper.sty}{\usepackage{footnotehyper}}{\usepackage{footnote}}
\makesavenoteenv{longtable}
\usepackage{graphicx}
\makeatletter
\newsavebox\pandoc@box
\newcommand*\pandocbounded[1]{% scales image to fit in text height/width
  \sbox\pandoc@box{#1}%
  \Gscale@div\@tempa{\textheight}{\dimexpr\ht\pandoc@box+\dp\pandoc@box\relax}%
  \Gscale@div\@tempb{\linewidth}{\wd\pandoc@box}%
  \ifdim\@tempb\p@<\@tempa\p@\let\@tempa\@tempb\fi% select the smaller of both
  \ifdim\@tempa\p@<\p@\scalebox{\@tempa}{\usebox\pandoc@box}%
  \else\usebox{\pandoc@box}%
  \fi%
}
% Set default figure placement to htbp
\def\fps@figure{htbp}
\makeatother
\setlength{\emergencystretch}{3em} % prevent overfull lines
\providecommand{\tightlist}{%
  \setlength{\itemsep}{0pt}\setlength{\parskip}{0pt}}
\usepackage{bookmark}
\IfFileExists{xurl.sty}{\usepackage{xurl}}{} % add URL line breaks if available
\urlstyle{same}
\hypersetup{
  hidelinks,
  pdfcreator={LaTeX via pandoc}}

\author{}
\date{}

\begin{document}

\section{HELIX: Diagnosing and Provably Restoring Post-Hoc Control Layer
Collapse via Threshold
Decoupling}\label{helix-diagnosing-and-provably-restoring-post-hoc-control-layer-collapse-via-threshold-decoupling}

\textbf{K Dhiraj}\\
Department of Computer Science and Engineering\\
Chandigarh University\\
k.dhiraj.srihari@gmail.com

\textbf{Abstract---} Post-hoc control layers impose inference-time
constraints on trained classifiers without modifying model weights.
HELIX is a margin-based module for class-selective override in a neural
IDS, designed to suppress low-confidence class-4 predictions. We
formally define \textbf{control collapse} --- a structural failure mode
in which the override layer processes every prediction but modifies none
--- using a dual criterion: override\_rate \textless{} \(\varepsilon\)
and KL divergence between pre- and post-override prediction
distributions equal to zero.

We prove that the \emph{online\_hybrid} configuration satisfies both
criteria exactly on UNSW-NB15 (override\_rate = 0.0, KL = 0.0) and
NSL-KDD (override\_rate = 0.0, KL = 0.0), despite class-4 prediction
mass of 20.6\% and 42.0\% respectively. The root cause is adaptive
threshold coupling: the threshold is estimated from the same margin
stream it is meant to bound, and a dimensional mismatch
(\(\tau_a \sim O(10^4)\) vs.~\(z_{\text{norm}} \sim O(1)\)) causes the
gate to fail silently.

A minimal corrective intervention --- freezing the adaptive threshold at
calibration time, where
\(\tau_a^{\text{frozen}} = \mu_{\text{cal}} + 0.754\sigma_{\text{cal}} = 81{,}921\)
logit units --- breaks the coupling and restores non-zero override
actuation. Bootstrap validation (n = 1,000 resamples, 95\% CI) confirms
statistically stable restoration: UNSW override\_rate = 0.0947
{[}0.09125, 0.09836{]}; NSL override\_rate = 0.3784 {[}0.37230,
0.38459{]}. The collapsed and active configurations have non-overlapping
CIs, with zero-width variance under collapse confirming deterministic
architectural failure rather than sampling noise.

One transparency result merits disclosure: under current calibration,
the decoupled \emph{frozen\_z\_hybrid} configuration is numerically
identical to the simpler \emph{fixed\_tau\_only} baseline. The z-score
normalization channel adds architectural robustness but no additional
discriminative signal at these parameter settings.

\textbf{Index Terms---} intrusion detection system, post-hoc control,
confidence thresholding, control collapse, conformal calibration,
selective classification, edge security.

\subsection{I. INTRODUCTION}\label{i.-introduction}

Post-hoc inference control is a practical alternative to full retraining
when decision-level constraints must be imposed on a deployed classifier
{[}1{]}--{[}4{]}. The standard use case in intrusion detection is
conservative reclassification: when the model assigns uncertain top-1
predictions to a high-cost class, override the prediction at inference
time rather than retrain the model to fix it.

The natural confidence signal for such a criterion is the
\textbf{decision margin} --- the gap between the top-1 and second-best
logit scores. Margin-based thresholding is used in confidence
calibration {[}5{]}, selective classification {[}6{]}, and reject-option
design {[}7{]}. A common refinement is the \textbf{adaptive threshold}:
compute the override boundary from the running statistics of the
observed margin distribution, so the violation rate stays roughly stable
even as the input distribution shifts.

Under a specific implementation pattern, adaptive thresholding can
nullify the control layer it is intended to enforce.

\textbf{The central result:} When the threshold is estimated online from
the same stream over which violations are evaluated, and when the
threshold estimator and the comparison operand live in different units,
the override probability converges to zero by construction. The control
layer runs on every sample and changes none. We call this failure
\textbf{control collapse}, define it formally, prove it holds in the
HELIX system, identify the mechanism, and show a minimal fix.

The failure mode is not exposed by existing methods. Confidence
calibration {[}5, 8{]} operates in probability space with static
parameters. OOD detection {[}9, 10{]} applies fixed, pre-computed
thresholds. Conformal prediction {[}11, 12{]} provides coverage
guarantees precisely \emph{because} it freezes the nonconformity
threshold before inference --- exactly the safeguard that the online
hybrid configuration violates. Selective classification {[}6, 13{]}
analyzes static coverage constraints; online threshold drift is out of
scope. This paper fills that gap.

\textbf{Contributions:} 1. \textbf{Formal collapse definition} via a
dual criterion (override\_rate \textless{} \(\varepsilon\), KL
divergence = 0) with exact measured values. 2. \textbf{Collapse proof}
showing \emph{online\_hybrid} satisfies both necessary and sufficient
conditions on two datasets. 3. \textbf{Mechanistic explanation} ---
Proposition 1 --- establishing coupling-induced convergence of override
probability to zero. 4. \textbf{Minimal corrective intervention} ---
frozen threshold decoupling --- restoring non-zero actuation, with
bootstrap confirmation of stability. 5.
\textbf{Controllability--correctness analysis} documenting the trade-off
and its calibration context, including the transparency result that
\emph{frozen\_z\_hybrid} \(\equiv\) \emph{fixed\_tau\_only} under
current parameters.

\begin{center}\rule{0.5\linewidth}{0.5pt}\end{center}

\subsection{II. RELATED WORK}\label{ii.-related-work}

\subsubsection{A. Confidence
Calibration}\label{a.-confidence-calibration}

Guo et al.~{[}5{]} show that modern neural networks are systematically
miscalibrated and propose temperature scaling as a fix. Platt scaling
{[}14{]} and isotonic regression {[}15{]} address the same objective via
different parametric forms. All these methods work in probability space
with a temperature parameter that is static after calibration. None
analyzes the behavior of an override threshold that is continuously
re-estimated from the test stream. The control collapse failure mode is
invisible to calibration analysis because calibration is not in the
loop.

\subsubsection{B. Out-of-Distribution
Detection}\label{b.-out-of-distribution-detection}

Hendrycks and Gimpel {[}9{]} propose maximum softmax probability as a
baseline OOD score; Lee et al.~{[}10{]} use Mahalanobis distance in
feature space; ODIN {[}16{]} combines temperature scaling and input
perturbations. All three apply a \textbf{fixed, pre-computed threshold}
to the score --- the coupling problem does not arise. Collapse is
specific to the regime where the threshold co-evolves with the score
distribution during evaluation.

\subsubsection{C. Conformal Prediction and Selective
Classification}\label{c.-conformal-prediction-and-selective-classification}

Conformal prediction {[}11, 12{]} provides valid coverage guarantees
under exchangeability by computing the nonconformity threshold on a
held-out calibration set and \textbf{freezing it before inference}.
Selective classification {[}6, 13{]} similarly fixes coverage
constraints at deployment. Both frameworks are structurally correct for
the same reason the frozen HELIX variant works: threshold and evaluation
stream are decoupled. The present work documents what happens when that
safeguard is removed.

\subsubsection{D. Adaptive Thresholding in Streaming
Settings}\label{d.-adaptive-thresholding-in-streaming-settings}

Concept drift detection {[}17, 18{]} uses adaptive thresholds on the
label or classifier space, not on the confidence threshold applied to
the classifier's own output. The self-referential pattern --- threshold
adapts to the same signal it is meant to bound --- is specific to
post-hoc control overrides. This paper provides a formal analysis of
that case.

\textbf{Why existing work misses this.} All surveyed methods either use
fixed thresholds from held-out calibration, operate in a different
representation space, or do not model self-referential threshold--margin
feedback. The \emph{online\_hybrid} configuration violates all three
safeguards simultaneously.

\begin{center}\rule{0.5\linewidth}{0.5pt}\end{center}

\subsection{III. PROBLEM SETUP AND FORMAL
DEFINITIONS}\label{iii.-problem-setup-and-formal-definitions}

\subsubsection{A. Margin Definition}\label{a.-margin-definition}

Let \(\mathbf{z}(x) \in \mathbb{R}^K\) denote the pre-softmax logit
vector from classifier \(f\) on input \(x\). HELIX computes margins in
\textbf{logit space}:

\[m = z_{(1)} - z_{(2)}\]

where \(z_{(1)}\) and \(z_{(2)}\) are the top-1 and second-best logit
values. Logit-space margins are required for dimensional consistency
with threshold values \(\tau_f = 70{,}000\) and
\(\tau_a^{\text{frozen}} = 81{,}921\) (both in logit units). All
threshold comparisons, z-score normalizations, and distributional
analyses in this paper are in logit margin space.

\subsubsection{B. Calibration
Statistics}\label{b.-calibration-statistics}

The following parameters are estimated from the UNSW-NB15 calibration
split and frozen before test evaluation:

{\def\LTcaptype{none} % do not increment counter
\begin{longtable}[]{@{}lll@{}}
\toprule\noalign{}
Parameter & Value & Units \\
\midrule\noalign{}
\endhead
\bottomrule\noalign{}
\endlastfoot
\(\mu_{\text{cal}}\) & 67,137.24 & logit units \\
\(\sigma_{\text{cal}}\) & 19,616.67 & logit units \\
\(\tau_a^{\text{frozen}}\) & 81,921.39 & logit units \\
\(\tau_f\) & 70,000.00 & logit units \\
\end{longtable}
}

The frozen threshold's normalized position:

\[z_{\tau_a} = \frac{\tau_a^{\text{frozen}} - \mu_{\text{cal}}}{\sigma_{\text{cal}}} = \frac{81{,}921.39 - 67{,}137.24}{19{,}616.67} \approx 0.754\]

This places \(\tau_a^{\text{frozen}}\) at approximately the \textbf{77th
percentile} of the calibration margin distribution, converting HELIX
from a dynamic estimator into a deterministic constraint function. The
fixed threshold \(\tau_f = 70{,}000\) sits at
\(z_{\tau_f} \approx 0.146\), or roughly the \textbf{56th percentile}.
In principle, the OR gate covers class-4 samples from the 56th to the
77th percentile of calibration confidence --- though as Section V shows,
the two gates produce identical decisions under current
parameterization.

\subsubsection{C. The Override Rule}\label{c.-the-override-rule}

\[\hat{y} = \begin{cases} \arg\max_{k \neq 4} z_k(x) & \text{if } y = 4 \;\wedge\; \left( m < \tau_f \;\vee\; \frac{m - \mu_{\text{cal}}}{\sigma_{\text{cal}}} < \tau_a \right) \\ y & \text{otherwise} \end{cases}\]

where \(y = \arg\max_k z_k\) and all parameters are as defined in
Section III.B.

\textbf{Figure 1} shows the full inference-time pipeline, from raw
logits through margin computation to the dual-gate threshold check and
override decision.

\begin{figure}
\centering
\includegraphics[width=0.72\linewidth,height=\textheight,keepaspectratio,alt={Figure 1. HELIX inference-time control pipeline. The dual-gate threshold check (fixed: m \textless{} 70\{,\}000; adaptive: z\_\{ ext\{norm\}\} \textless{}    au\_a) evaluates an OR condition on every class-4 top prediction.}]{docs/figures/fig1.png}
\caption{Figure 1. HELIX inference-time control pipeline. The dual-gate
threshold check (fixed: \(m < 70{,}000\); adaptive:
\(z_{ ext{norm}} <    au_a\)) evaluates an OR condition on every class-4
top prediction.}
\end{figure}

The control layer is logically well-formed; inactivity under
\emph{online\_hybrid} is not a logic error --- it is gate failure from
null-threshold propagation, as established in Section IV.

\subsubsection{D. Ablation
Configurations}\label{d.-ablation-configurations}

{\def\LTcaptype{none} % do not increment counter
\begin{longtable}[]{@{}
  >{\raggedright\arraybackslash}p{(\linewidth - 8\tabcolsep) * \real{0.2000}}
  >{\raggedright\arraybackslash}p{(\linewidth - 8\tabcolsep) * \real{0.2000}}
  >{\raggedright\arraybackslash}p{(\linewidth - 8\tabcolsep) * \real{0.2000}}
  >{\raggedright\arraybackslash}p{(\linewidth - 8\tabcolsep) * \real{0.2000}}
  >{\raggedright\arraybackslash}p{(\linewidth - 8\tabcolsep) * \real{0.2000}}@{}}
\toprule\noalign{}
\begin{minipage}[b]{\linewidth}\raggedright
Configuration
\end{minipage} & \begin{minipage}[b]{\linewidth}\raggedright
\(\tau_f\)
\end{minipage} & \begin{minipage}[b]{\linewidth}\raggedright
\(\tau_a\)
\end{minipage} & \begin{minipage}[b]{\linewidth}\raggedright
\(\tau_a\) update
\end{minipage} & \begin{minipage}[b]{\linewidth}\raggedright
\(\tau_a\) (logit)
\end{minipage} \\
\midrule\noalign{}
\endhead
\bottomrule\noalign{}
\endlastfoot
\texttt{fixed\_tau\_only} & 70,000 & --- & N/A & --- \\
\texttt{percentile\_only} & 0 & online & Running percentile & null at
eval \\
\texttt{online\_hybrid} & 70,000 & online & Running percentile & null at
eval \\
\texttt{frozen\_z\_hybrid} & 70,000 & frozen & Calibration &
81,921.39 \\
\end{longtable}
}

\subsubsection{E. Algorithm 1: HELIX Inference Control
(frozen\_z\_hybrid)}\label{e.-algorithm-1-helix-inference-control-frozen_z_hybrid}

\begin{verbatim}
Algorithm 1: HELIX Inference-Time Override (frozen_z_hybrid variant)
-------------------------------------------------------------------
Input:  z ∈ ℝ^K          — raw logit vector
        τ_f = 70,000      — fixed logit-margin threshold
        τ_a = 81,921.39   — frozen adaptive threshold (≈ 77th pctile of cal.)
        μ_cal = 67,137.24, σ_cal = 19,616.67  — calibration statistics

Output: ŷ ∈ {0,...,K−1}

 1:  y      ← argmax_k z_k                    // O(K)
 2:  if y ≠ 4 then return y                   // no override candidate
 3:  z_(1) ← max(z);  z_(2) ← second_max(z)  // O(K)
 4:  m      ← z_(1) - z_(2)                   // O(1)
 5:  z_norm ← (m - μ_cal) / σ_cal             // O(1)
 6:  gate_f ← (m < τ_f)                       // O(1)
 7:  gate_a ← (z_norm < τ_a)                  // O(1) — τ_a is a scalar constant
 8:  if gate_f OR gate_a then
 9:      return argmax_{k≠4} z_k              // O(K), override class 4
10:  else
11:      return 4
-------------------------------------------------------------------
\end{verbatim}

\subsubsection{F. Computational
Overhead}\label{f.-computational-overhead}

HELIX adds strictly \textbf{O(1)} computation per sample beyond the base
model forward pass:

{\def\LTcaptype{none} % do not increment counter
\begin{longtable}[]{@{}lll@{}}
\toprule\noalign{}
Step & Complexity & Notes \\
\midrule\noalign{}
\endhead
\bottomrule\noalign{}
\endlastfoot
argmax / top-2 & O(K) & Unavoidable; same as base model \\
Margin, z-normalization & O(1) & 3 scalar arithmetic operations \\
Gate evaluation & O(1) & 2 comparisons + 1 OR \\
Override argmax & O(K) & Over K−1 classes \\
\textbf{HELIX overhead} & \textbf{O(1)} & All post-argmax operations \\
\end{longtable}
}

Memory overhead is four scalar constants
(\(\tau_f, \tau_a^{\text{frozen}}, \mu_{\text{cal}}, \sigma_{\text{cal}}\)),
independent of model size, batch size, or sequence length. No
retraining, no gradient computation, no additional layers.

\begin{center}\rule{0.5\linewidth}{0.5pt}\end{center}

\subsection{IV. FAILURE MECHANISM AND
DIAGNOSIS}\label{iv.-failure-mechanism-and-diagnosis}

\subsubsection{A. Formal Definition}\label{a.-formal-definition}

\textbf{Definition 1 (Control Collapse --- Dual Criterion).} Let
\(\mathcal{M} = \{x : y(x) = 4\}\) be the set of class-4 candidate
samples, and let \(P(y)\) and \(P(\hat{y})\) denote the prediction
distributions before and after HELIX override. A HELIX control layer is
in \textbf{control collapse} if and only if both hold:

\textbf{Criterion 1 (Actuation failure):}
\[r = \frac{|\{x \in \mathcal{D} : \text{override}(x)\}|}{|\mathcal{D}|} < \varepsilon, \quad \varepsilon = 0.001\]

\textbf{Criterion 2 (Distribution invariance):}
\[D_{\text{KL}}\!\left(P(\hat{y}) \;\|\; P(y)\right) = 0\]

The dual criterion is necessary because Criterion 1 alone could be
satisfied when class-4 candidate mass is genuinely low, while Criterion
2 alone could hold trivially. Together they establish that the control
layer produces zero actuation \emph{and} leaves the prediction
distribution unchanged --- it is a null operator. Note that
\(D_{\text{KL}} = 0\) establishes distributional identity, while
\(D_{\text{KL}} > 10\) (as seen in active configurations) indicates
strong distributional intervention rather than marginal perturbation.

\subsubsection{B. Collapse Proof for
online\_hybrid}\label{b.-collapse-proof-for-online_hybrid}

\textbf{Claim.} \emph{online\_hybrid} is in control collapse on both
UNSW-NB15 and NSL-KDD.

\textbf{Proof by direct measurement:}

\emph{UNSW-NB15 (n = 26,302 samples):} - Criterion 1: override\_rate =
0.0 \textless{} 0.001 (satisfied; override\_count = 0) - Criterion 2:
\(D_{\text{KL}}(P(\hat{y}) \| P(y)) = 0.0\) (satisfied) -
\texttt{collapsed} = \textbf{True}

\emph{NSL-KDD (n = 22,278 samples):} - Criterion 1: override\_rate = 0.0
\textless{} 0.001 (satisfied; override\_count = 0) - Criterion 2:
\(D_{\text{KL}}(P(\hat{y}) \| P(y)) = 0.0\) (satisfied) -
\texttt{collapsed} = \textbf{True}

Both criteria hold exactly --- \(D_{\text{KL}} = 0.0\) is not numerical
roundoff. It reflects bit-identical pre- and post-override prediction
vectors, confirmed by metric identity between \emph{online\_hybrid} and
\emph{percentile\_only} across all performance measures. \textbf{Figure
3 (left panels)} renders this distributional identity visually: the pre-
and post-override margin histograms are indistinguishable.

\textbf{Thus, \emph{online\_hybrid} satisfies both the necessary and
sufficient conditions for control collapse across both datasets.}
\(\square\)

\textbf{Table 3: Collapse quantification}

{\def\LTcaptype{none} % do not increment counter
\begin{longtable}[]{@{}
  >{\raggedright\arraybackslash}p{(\linewidth - 8\tabcolsep) * \real{0.2000}}
  >{\raggedright\arraybackslash}p{(\linewidth - 8\tabcolsep) * \real{0.2000}}
  >{\raggedright\arraybackslash}p{(\linewidth - 8\tabcolsep) * \real{0.2000}}
  >{\raggedright\arraybackslash}p{(\linewidth - 8\tabcolsep) * \real{0.2000}}
  >{\raggedright\arraybackslash}p{(\linewidth - 8\tabcolsep) * \real{0.2000}}@{}}
\toprule\noalign{}
\begin{minipage}[b]{\linewidth}\raggedright
Dataset
\end{minipage} & \begin{minipage}[b]{\linewidth}\raggedright
Configuration
\end{minipage} & \begin{minipage}[b]{\linewidth}\raggedright
Override rate
\end{minipage} & \begin{minipage}[b]{\linewidth}\raggedright
KL divergence
\end{minipage} & \begin{minipage}[b]{\linewidth}\raggedright
Collapsed state
\end{minipage} \\
\midrule\noalign{}
\endhead
\bottomrule\noalign{}
\endlastfoot
UNSW & \texttt{fixed\_tau\_only} & 0.0947 & 10.218 & \textbf{No} \\
UNSW & \texttt{percentile\_only} & 0.0000 & 0.000 & Yes \\
UNSW & \texttt{online\_hybrid} & \textbf{0.0000} & \textbf{0.000} &
Yes \\
UNSW & \texttt{frozen\_z\_hybrid} & 0.0947 & 10.218 & \textbf{No} \\
NSL & \texttt{fixed\_tau\_only} & 0.3784 & 22.922 & \textbf{No} \\
NSL & \texttt{percentile\_only} & 0.0000 & 0.000 & Yes \\
NSL & \texttt{online\_hybrid} & \textbf{0.0000} & \textbf{0.000} &
Yes \\
NSL & \texttt{frozen\_z\_hybrid} & 0.3784 & 22.922 & \textbf{No} \\
\end{longtable}
}

\emph{Collapse threshold \(\varepsilon\) = 0.001. KL metric:
\(D_{\text{KL}}(\text{pre\_top1} \| \text{post\_pred})\) over class-4
margin distributions.}

The KL values of 10.218 (UNSW) and 22.922 (NSL) for active
configurations show that override intervention produces substantial
redistribution of the class-4 margin density --- not a marginal shift.
This contrast with the collapsed configurations
(\(D_{\text{KL}} = 0.0\)) is the core empirical result.

\subsubsection{C. Mechanism: Coupling-Induced
Collapse}\label{c.-mechanism-coupling-induced-collapse}

\textbf{Proposition 1 (Coupling-Induced Override Collapse).} \emph{Let
\(\{m_t\}\) be a bounded sequence of logit-space margins. Let
\(\tau_a(t) = \hat{F}^{-1}_M(q, t)\) be the \(q\)-th empirical quantile
of \(\{m_1, \ldots, m_t\}\) (online estimate). Suppose the adaptive gate
evaluates \(z_{\text{norm}}(m_t) < \tau_a(t)\), where
\(z_{\text{norm}}\) is dimensionless and \(\tau_a(t)\) is in logit
units. If \(\tau_a(t) \gg \max_t z_{\text{norm}}(m_t)\) (unit mismatch
regime), the gate is trivially satisfied, reducing the effective
override decision to the fixed gate alone. If the fixed gate is also
disabled (\(\tau_f = \text{null}\) or \(\tau_f = 0\)), then
\(P(\text{override}) = 0\) exactly.}

\textbf{Proof sketch.} In \emph{online\_hybrid}, \texttt{tau\_adaptive}
is recorded as \texttt{null} at evaluation time (confirmed by
\texttt{"tau\_adaptive":\ null} in the experimental record). A
null-valued threshold in the comparison
\texttt{z\_norm\ \textless{}\ null} evaluates to False in the runtime.
The fixed gate \(m < \tau_f = 70{,}000\) is present, but a bug (noted in
\texttt{bug\_fixes}: ``fixed gate uses raw tau positivity'') disables it
when the adaptive gate is null. Both gates fail silently, producing
override\_count = 0 for all inputs. In \emph{percentile\_only},
\(\tau_f = 0\) renders \(m < 0\) false for all non-negative margins.
Both failure paths reach the same outcome. \(\square\)

\textbf{Corollary.} \emph{online\_hybrid} \(\equiv\)
\emph{percentile\_only} exactly: both reduce to a no-op override layer
through structurally different but mutually reinforcing gate failures.

\subsubsection{D. Why This Failure Is
Structural}\label{d.-why-this-failure-is-structural}

\[\underbrace{\text{Online coupling}}_{\tau_a(t) \propto m_t} \;\times\; \underbrace{\text{Unit mismatch}}_{\tau_a \sim O(10^4),\; z_{\text{norm}} \sim O(1)} \;\Rightarrow\; \underbrace{\text{Gate null or trivial}}_{\text{no discriminative comparison}} \;\Rightarrow\; \underbrace{r = 0, \; D_{\text{KL}} = 0}_{\text{collapse}}\]

This is a property of the coupled online estimation architecture, not of
the specific data or calibration values.

\textbf{Figure 2} shows the causal chain side by side: the collapse path
(left, red) and the restoration path (right, green).

\begin{figure}
\centering
\includegraphics[width=0.72\linewidth,height=\textheight,keepaspectratio,alt={Figure 2. Causal mechanism of control collapse (left) and restoration by frozen-threshold decoupling (right).}]{docs/figures/fig2.png}
\caption{Figure 2. Causal mechanism of control collapse (left) and
restoration by frozen-threshold decoupling (right).}
\end{figure}

Left: online percentile estimation produces a null \$ au\_a\$ at
evaluation; the gate fails silently and \(D_{   ext{KL}} = 0\). Right:
freezing \$ au\_a\^{}\{ ext\{frozen\}\}\$ at calibration provides a
stable reference; gates activate and \(D_{   ext{KL}} > 10\).

\begin{center}\rule{0.5\linewidth}{0.5pt}\end{center}

\subsection{V. INTERVENTION AND SCOPE OF
RESTORATION}\label{v.-intervention-and-scope-of-restoration}

\subsubsection{A. Principle}\label{a.-principle}

Freezing \(\tau_a\) at calibration eliminates the feedback loop. The
threshold becomes a scalar constant computed once (from the UNSW
calibration split: \(\tau_a^{\text{frozen}} = 81{,}921.39\)) and applied
unchanged at inference. As established in Section III.B, this placement
at \(\mu_{\text{cal}} + 0.754\sigma_{\text{cal}}\) defines a fixed
quantile operator over the calibration distribution --- a deterministic
constraint rather than a running estimate.

\textbf{Figure 4 (right panel)} shows the decoupled threshold sitting
fixed at 81,921 while the evaluation batch margins move freely above and
below it, enabling violation detection. The left panel shows the coupled
case: \(\tau_{\text{adaptive}}\) chases the distribution, and the gap
between threshold and mean margin collapses.

\begin{figure}
\centering
\includegraphics[width=0.72\linewidth,height=\textheight,keepaspectratio,alt={Figure 4. Threshold and class-4 margin evolution across evaluation batches for coupled (online\_hybrid) and decoupled (frozen\_z\_hybrid) operation.}]{docs/figures/fig4.png}
\caption{Figure 4. Threshold and class-4 margin evolution across
evaluation batches for coupled (\emph{online\_hybrid}) and decoupled
(\emph{frozen\_z\_hybrid}) operation.}
\end{figure}

In the coupled case, the adaptive threshold tracks the mean margin. In
the decoupled case, \$ au\_a\^{}\{ ext\{frozen\}\} = 81\{,\}921\$
remains fixed, allowing genuine margin-threshold crossings.

\subsubsection{B. Controllability Restored:
Results}\label{b.-controllability-restored-results}

\textbf{Figure 5} shows per-batch override counts across both methods
and datasets. The left panel (\emph{online\_hybrid}) is uniformly zero.
The right (\emph{frozen\_z\_hybrid}) shows consistent non-zero
activation across all batches, totaling 2,491 overrides on UNSW and
8,430 on NSL.

\begin{figure}
\centering
\includegraphics[width=0.72\linewidth,height=\textheight,keepaspectratio,alt={Figure 5. Per-batch override activation counts for online\_hybrid and frozen\_z\_hybrid on UNSW and NSL.}]{docs/figures/fig5.png}
\caption{Figure 5. Per-batch override activation counts for
\emph{online\_hybrid} and \emph{frozen\_z\_hybrid} on UNSW and NSL.}
\end{figure}

\emph{online\_hybrid} produces zero overrides in every batch on both
datasets, while \emph{frozen\_z\_hybrid} exhibits stable non-zero
activation.

\textbf{Figure 3 (right panels)} shows the distributional consequence:
the post-override class-4 margin distribution is rightward-shifted in
the active configuration (\(D_{\text{KL}} = 10.218\) on UNSW,
\(D_{\text{KL}} = 22.922\) on NSL).

\begin{figure}
\centering
\includegraphics[width=0.84\linewidth,height=\textheight,keepaspectratio,alt={Figure 3. Class-4 logit margin distributions (pre vs.~post override) for collapsed and active regimes on UNSW and NSL.}]{docs/figures/fig3.png}
\caption{Figure 3. Class-4 logit margin distributions (pre vs.~post
override) for collapsed and active regimes on UNSW and NSL.}
\end{figure}

For \emph{online\_hybrid}, pre and post distributions are identical
(\(D_{ ext{KL}} = 0\)). For \emph{frozen\_z\_hybrid}, post distributions
shift substantially (\(D_{ ext{KL}} = 10.218\) on UNSW, 22.922 on NSL).

\subsubsection{\texorpdfstring{C. frozen\_z\_hybrid \(\equiv\)
fixed\_tau\_only: A Transparency
Note}{C. frozen\_z\_hybrid \textbackslash equiv fixed\_tau\_only: A Transparency Note}}\label{c.-frozen_z_hybrid-equiv-fixed_tau_only-a-transparency-note}

Under current calibration, \emph{frozen\_z\_hybrid} produces identical
override decisions to \emph{fixed\_tau\_only} for every sample in both
datasets. The z-score gate
(\(z_{\text{norm}} < \tau_a^{\text{frozen}}\)) and the fixed gate
(\(m < \tau_f\)) select the same overrides --- the lower-percentile gate
(56th) subsumes the higher-percentile gate (77th) in all observed cases,
collapsing the OR condition to a single effective criterion.

This means normalization provides architectural robustness --- if
\(\sigma_{\text{cal}}\) changes substantially, the relative threshold
placement is preserved --- but does not add discriminative signal at
present. Recalibrating \(\tau_a^{\text{frozen}}\) to a percentile
distinct from \(\tau_f\) would expose independent z-score gate behavior.

\begin{center}\rule{0.5\linewidth}{0.5pt}\end{center}

\subsection{VI. EXPERIMENTAL PROTOCOL}\label{vi.-experimental-protocol}

\subsubsection{A. Datasets}\label{a.-datasets}

\textbf{UNSW-NB15} {[}19{]}: Contemporary network intrusion detection
dataset. The base model is trained on the UNSW-NB15 training partition;
evaluation on the test partition is the \textbf{in-distribution}
protocol (26,302 samples).

\textbf{NSL-KDD} {[}20{]}: Curated version of KDD Cup 1999 with
corrected statistical biases. Feature distribution differs substantially
from UNSW-NB15. Evaluating the UNSW-trained checkpoint on NSL-KDD is the
\textbf{cross-distribution} protocol (22,278 samples). No fine-tuning or
adaptation is performed.

\subsubsection{B. Bootstrap Validation}\label{b.-bootstrap-validation}

Statistical confidence intervals are computed by non-parametric
bootstrap: - Resamples: 1,000 (with replacement) - Confidence level:
95\% - Strategy: resample prediction-level outcomes (override decisions
and class assignments)

Exact bootstrap CIs are available for override\_rate (sample-level
metric). Macro-F1, class4\_precision, and class4\_recall require
sample-level prediction vectors; current data export provides
histogram-level approximations only, so those CIs are marked
\texttt{approx\_from\_histogram} in the experimental record and are not
reported here.

\begin{center}\rule{0.5\linewidth}{0.5pt}\end{center}

\subsection{VII. RESULTS}\label{vii.-results}

\subsubsection{A. Performance
Comparison}\label{a.-performance-comparison}

\textbf{Table 1: Classification performance across configurations}

{\def\LTcaptype{none} % do not increment counter
\begin{longtable}[]{@{}
  >{\raggedright\arraybackslash}p{(\linewidth - 10\tabcolsep) * \real{0.1667}}
  >{\raggedright\arraybackslash}p{(\linewidth - 10\tabcolsep) * \real{0.1667}}
  >{\raggedright\arraybackslash}p{(\linewidth - 10\tabcolsep) * \real{0.1667}}
  >{\raggedright\arraybackslash}p{(\linewidth - 10\tabcolsep) * \real{0.1667}}
  >{\raggedright\arraybackslash}p{(\linewidth - 10\tabcolsep) * \real{0.1667}}
  >{\raggedright\arraybackslash}p{(\linewidth - 10\tabcolsep) * \real{0.1667}}@{}}
\toprule\noalign{}
\begin{minipage}[b]{\linewidth}\raggedright
Dataset
\end{minipage} & \begin{minipage}[b]{\linewidth}\raggedright
Configuration
\end{minipage} & \begin{minipage}[b]{\linewidth}\raggedright
Macro-F1
\end{minipage} & \begin{minipage}[b]{\linewidth}\raggedright
Class-4 precision
\end{minipage} & \begin{minipage}[b]{\linewidth}\raggedright
Class-4 recall
\end{minipage} & \begin{minipage}[b]{\linewidth}\raggedright
Class-4 prediction ratio
\end{minipage} \\
\midrule\noalign{}
\endhead
\bottomrule\noalign{}
\endlastfoot
UNSW & \texttt{fixed\_tau\_only} & \textbf{0.3411} & \textbf{0.0563} &
\textbf{0.9706} & 0.1115 \\
UNSW & \texttt{percentile\_only} & 0.2626 & 0.0313 & 1.0000 & 0.2062 \\
UNSW & \texttt{online\_hybrid} & 0.2626 & 0.0313 & 1.0000 & 0.2062 \\
UNSW & \texttt{frozen\_z\_hybrid} & \textbf{0.3411} & \textbf{0.0563} &
\textbf{0.9706} & 0.1115 \\
NSL & \texttt{fixed\_tau\_only} & \textbf{0.0226} & \textbf{0.00108} &
\textbf{0.0263} & 0.0416 \\
NSL & \texttt{percentile\_only} & 0.0220 & 0.000321 & 0.0789 & 0.4200 \\
NSL & \texttt{online\_hybrid} & 0.0220 & 0.000321 & 0.0789 & 0.4200 \\
NSL & \texttt{frozen\_z\_hybrid} & \textbf{0.0226} & \textbf{0.00108} &
\textbf{0.0263} & 0.0416 \\
\end{longtable}
}

\emph{Note: Macro-F1 values are deterministic aggregates without
bootstrap CIs (see Section VI.B). The metric identity between}
online\_hybrid \emph{and} percentile\_only \emph{across all columns
confirms control collapse at the performance level.}

\subsubsection{B. Override Actuation and Confidence
Intervals}\label{b.-override-actuation-and-confidence-intervals}

\textbf{Table 2: Control activation with bootstrap confidence intervals}

{\def\LTcaptype{none} % do not increment counter
\begin{longtable}[]{@{}
  >{\raggedright\arraybackslash}p{(\linewidth - 10\tabcolsep) * \real{0.1667}}
  >{\raggedright\arraybackslash}p{(\linewidth - 10\tabcolsep) * \real{0.1667}}
  >{\raggedright\arraybackslash}p{(\linewidth - 10\tabcolsep) * \real{0.1667}}
  >{\raggedright\arraybackslash}p{(\linewidth - 10\tabcolsep) * \real{0.1667}}
  >{\raggedright\arraybackslash}p{(\linewidth - 10\tabcolsep) * \real{0.1667}}
  >{\raggedright\arraybackslash}p{(\linewidth - 10\tabcolsep) * \real{0.1667}}@{}}
\toprule\noalign{}
\begin{minipage}[b]{\linewidth}\raggedright
Dataset
\end{minipage} & \begin{minipage}[b]{\linewidth}\raggedright
Configuration
\end{minipage} & \begin{minipage}[b]{\linewidth}\raggedright
Override count
\end{minipage} & \begin{minipage}[b]{\linewidth}\raggedright
Override rate
\end{minipage} & \begin{minipage}[b]{\linewidth}\raggedright
95\% CI
\end{minipage} & \begin{minipage}[b]{\linewidth}\raggedright
\(\sigma_{\text{override}}\)
\end{minipage} \\
\midrule\noalign{}
\endhead
\bottomrule\noalign{}
\endlastfoot
UNSW & \texttt{fixed\_tau\_only} & 2,491 & 0.0947 & --- & --- \\
UNSW & \texttt{percentile\_only} & 0 & 0.0000 & --- & --- \\
UNSW & \texttt{online\_hybrid} & \textbf{0} & \textbf{0.0000} &
\textbf{{[}0.000, 0.000{]}} & \textbf{0.00000} \\
UNSW & \texttt{frozen\_z\_hybrid} & 2,491 & \textbf{0.0947} &
\textbf{{[}0.091, 0.098{]}} & \textbf{0.00182} \\
NSL & \texttt{fixed\_tau\_only} & 8,430 & 0.3784 & --- & --- \\
NSL & \texttt{percentile\_only} & 0 & 0.0000 & --- & --- \\
NSL & \texttt{online\_hybrid} & \textbf{0} & \textbf{0.0000} &
\textbf{{[}0.000, 0.000{]}} & \textbf{0.00000} \\
NSL & \texttt{frozen\_z\_hybrid} & 8,430 & \textbf{0.3784} &
\textbf{{[}0.372, 0.385{]}} & \textbf{0.00314} \\
\end{longtable}
}

\emph{Bootstrap CIs from 1,000 resamples with replacement, 95\%
confidence level. fixed\_tau\_only and percentile\_only CIs are omitted
as they are structurally identical to frozen\_z\_hybrid and
online\_hybrid respectively.}

\textbf{Figure 6} plots the override rate CIs for all four evaluated
configurations.

\begin{figure}
\centering
\includegraphics[width=0.78\linewidth,height=\textheight,keepaspectratio,alt={Figure 6. Bootstrap confidence intervals for override rate (95\%, n=1,000) across collapsed and active configurations.}]{docs/figures/fig6.png}
\caption{Figure 6. Bootstrap confidence intervals for override rate
(95\%, n=1,000) across collapsed and active configurations.}
\end{figure}

Collapsed configurations map to a point mass at zero (zero width, zero
variance). Active configurations have strictly non-overlapping non-zero
intervals.

\subsubsection{C. Statistical
Interpretation}\label{c.-statistical-interpretation}

\textbf{Non-overlapping CIs establish significant separation:}

\begin{itemize}
\tightlist
\item
  UNSW: \emph{frozen\_z\_hybrid} CI = {[}0.09125, 0.09836{]}
  vs.~\emph{online\_hybrid} CI = {[}0.000, 0.000{]}. Gap: 0.09125
  between the lower bound of active and the upper bound of collapsed.
\item
  NSL: \emph{frozen\_z\_hybrid} CI = {[}0.37230, 0.38459{]}
  vs.~\emph{online\_hybrid} CI = {[}0.000, 0.000{]}. Gap: 0.37230.
\end{itemize}

Both separations are statistically significant at 95\% confidence.

\textbf{Degenerate variance is a structural result, not a data
artifact.} The bootstrap standard deviation for \emph{online\_hybrid}
override\_rate is \(\sigma_{\text{override}} = 0.00000\) on both
datasets. Zero variance under 1,000 bootstrap resamples means no
resampling of the evaluation data can produce non-zero actuation. This
directly implies the failure is a property of the control architecture
and threshold state --- not of any particular subset of the test data.
Combined with \(D_{\text{KL}} = 0.0\), the collapse is established at
both the distributional and sample levels.

\begin{center}\rule{0.5\linewidth}{0.5pt}\end{center}

\subsection{VIII. DISCUSSION}\label{viii.-discussion}

\subsubsection{A. Controllability and Classification
Trade-off}\label{a.-controllability-and-classification-trade-off}

On UNSW, active control raises macro\_f1 from 0.2626 to 0.3411 (+0.0785)
by suppressing high-volume, low-precision class-4 predictions. Class-4
precision increases from 0.0313 to 0.0563 at a cost: recall drops from
1.000 to 0.9706. The mechanism targets samples below the 56th
calibration percentile of margin (the fixed gate), suppressing the
lowest-confidence class-4 predictions. The KL divergence of 10.218
establishes that the intervention is substantial --- not a borderline
perturbation.

\subsubsection{B. Cross-Distribution
Results}\label{b.-cross-distribution-results}

NSL macro\_f1 \textless{} 0.023 across all configurations is not a HELIX
failure --- it reflects representation collapse under distribution
shift. The model was trained on UNSW; it does not encode NSL class
structure in its logit space, and no post-hoc control layer can fix
that.

What the NSL results do show is a miscalibration problem. The 37.84\%
override rate on NSL (95\% CI: {[}37.23\%, 38.46\%{]}) means the
UNSW-calibrated threshold fires on more than a third of the NSL test
set. NSL class-4 margins are systematically lower under the UNSW model
--- the model is uncertain, and the uncertainty is genuine, not a
control artifact. The high KL divergence (22.922) reflects a
large-volume intervention, not a targeted one.

\subsubsection{C. What the Degenerate Variance Tells
Us}\label{c.-what-the-degenerate-variance-tells-us}

The \(\sigma_{\text{override}} = 0.00000\) result for
\emph{online\_hybrid} is more informative than a point estimate. It
establishes \textbf{sample-order invariance}: collapse holds regardless
of which subset of the evaluation data is examined. This rules out
explanations based on aggregation, imbalanced batches, or sampling
noise. Collapse is a function of the architecture and threshold state,
not of the data.

\begin{center}\rule{0.5\linewidth}{0.5pt}\end{center}

\subsubsection{D. Control-theoretic
interpretation}\label{d.-control-theoretic-interpretation}

HELIX maps onto discrete-time control system semantics:

{\def\LTcaptype{none} % do not increment counter
\begin{longtable}[]{@{}
  >{\raggedright\arraybackslash}p{(\linewidth - 4\tabcolsep) * \real{0.3333}}
  >{\raggedright\arraybackslash}p{(\linewidth - 4\tabcolsep) * \real{0.3333}}
  >{\raggedright\arraybackslash}p{(\linewidth - 4\tabcolsep) * \real{0.3333}}@{}}
\toprule\noalign{}
\begin{minipage}[b]{\linewidth}\raggedright
Control concept
\end{minipage} & \begin{minipage}[b]{\linewidth}\raggedright
HELIX equivalent
\end{minipage} & \begin{minipage}[b]{\linewidth}\raggedright
Value
\end{minipage} \\
\midrule\noalign{}
\endhead
\bottomrule\noalign{}
\endlastfoot
State signal & Logit margin \(m_t\) & \textasciitilde67,137 ± 19,617
(calibration) \\
Control boundary & Override threshold \(\tau\) & \(\tau_f\) = 70,000;
\(\tau_a^{\text{frozen}}\) = 81,921 \\
Actuation & Override decision & 0\% (collapsed) vs.~9.5\%--37.8\%
(active) \\
Controller & Dual-gate evaluator & O(1) per sample \\
Plant & Base classifier \(f\) & helix\_full\_unsw\_nb15\_best.pt \\
\end{longtable}
}

Under \emph{online\_hybrid}, the controller redefines the boundary to
contain all observed states --- equivalent to a setpoint that chases the
process variable --- producing zero actuation authority. Under
\emph{frozen\_z\_hybrid}, the boundary is set externally at calibration
and held fixed, so the margin can genuinely cross it.

The KL divergence of 10.218 (UNSW) and 22.922 (NSL) under active
configurations quantifies the resulting actuation effect on the
prediction distribution. These are not marginal adjustments.

\begin{center}\rule{0.5\linewidth}{0.5pt}\end{center}

\subsection{IX. OPERATIONAL COMPATIBILITY AND DEPLOYMENT
CONSTRAINTS}\label{ix.-operational-compatibility-and-deployment-constraints}

Collapse resolution is not the same as deployment approval. We define
the \textbf{bounded control regime} as:

\[0 < \text{override\_rate} \leq 0.02\]

This captures non-zero actuation authority while bounding intervention
intensity to an operationally safe envelope.

{\def\LTcaptype{none} % do not increment counter
\begin{longtable}[]{@{}
  >{\raggedright\arraybackslash}p{(\linewidth - 4\tabcolsep) * \real{0.3333}}
  >{\raggedright\arraybackslash}p{(\linewidth - 4\tabcolsep) * \real{0.3333}}
  >{\raggedright\arraybackslash}p{(\linewidth - 4\tabcolsep) * \real{0.3333}}@{}}
\toprule\noalign{}
\begin{minipage}[b]{\linewidth}\raggedright
HELIX state
\end{minipage} & \begin{minipage}[b]{\linewidth}\raggedright
Control property
\end{minipage} & \begin{minipage}[b]{\linewidth}\raggedright
Deployment outcome
\end{minipage} \\
\midrule\noalign{}
\endhead
\bottomrule\noalign{}
\endlastfoot
\texttt{collapsed} & Inactive (override\_rate = 0) & PASS --- safe but
inactive \\
\texttt{restored} & Controllable (override\_rate \textgreater{} 0) &
FAIL under runbook constraints \\
\end{longtable}
}

The restored configurations (UNSW: 9.47\%, NSL: 37.84\%) exceed this
envelope. Restoration demonstrates that the control layer \emph{can}
act; it does not guarantee that it acts within operationally acceptable
bounds. That requires recalibration targeting the bounded regime, which
is outside the scope of the present paper.

\begin{center}\rule{0.5\linewidth}{0.5pt}\end{center}

\subsection{X. LIMITATIONS}\label{x.-limitations}

\begin{enumerate}
\def\labelenumi{\arabic{enumi}.}
\item
  \textbf{No true class signal recovery.} HELIX overrides argmax
  predictions but cannot improve underlying logit quality. Under
  representation collapse, the next-best class is equally unreliable.
\item
  \textbf{Cross-distribution miscalibration.} The 37.84\% NSL override
  rate is operationally pathological. Thresholds calibrated on UNSW do
  not produce bounded control on NSL.
\item
  \textbf{Class-4 scope only.} Override logic targets class 4
  exclusively. Multi-class extension requires per-class thresholds or a
  generalized margin criterion.
\item
  \textbf{frozen\_z\_hybrid \(\equiv\) fixed\_tau\_only under current
  calibration.} The z-score normalization channel provides no additional
  discriminative signal. Any benefit is fully attributable to the
  fixed-threshold gate. This is disclosed without mitigation.
\item
  \textbf{Macro-F1 bootstrap CIs unavailable.} Reported macro-F1 values
  are deterministic aggregates and should not be treated as
  statistically validated estimates.
\item
  \textbf{Single model, single training run.} Cross-model generalization
  of the collapse mechanism is established by Proposition 1
  (architectural argument), not empirical replication across models.
\end{enumerate}

\begin{center}\rule{0.5\linewidth}{0.5pt}\end{center}

\subsection{XI. CONCLUSION}\label{xi.-conclusion}

We formally defined control collapse via a dual criterion
(override\_rate \textless{} \(\varepsilon\), \(D_{\text{KL}} = 0\)) and
proved that \emph{online\_hybrid} satisfies both necessary and
sufficient conditions exactly on UNSW-NB15 (override\_rate = 0.0, KL =
0.0) and NSL-KDD (override\_rate = 0.0, KL = 0.0). The mechanism is
adaptive threshold coupling combined with null-value propagation through
the gate evaluation path. Freezing the threshold at calibration breaks
both failure modes simultaneously.

Bootstrap validation (1,000 resamples, 95\% CI) confirms that the
decoupled \emph{frozen\_z\_hybrid} configuration restores non-zero
actuation with non-overlapping CIs relative to the collapsed regime:
UNSW {[}0.091, 0.098{]} vs.~{[}0.000, 0.000{]}; NSL {[}0.372, 0.385{]}
vs.~{[}0.000, 0.000{]}. Zero-width variance under collapse confirms the
failure is deterministic and sample-order invariant.

One practical consequence extends beyond the HELIX system: \textbf{any
post-hoc control layer that re-estimates its violation threshold from
the live test stream is at structural risk of collapse under unit
mismatch.} Controllability --- a necessary prerequisite for correctness
--- cannot coexist with self-referential threshold estimation. In this
setting, freezing the threshold is a design requirement rather than a
heuristic workaround. Conformal prediction frameworks have known this
for decades, which is why they mandate held-out calibration. The lesson
from this failure mode is that the mandate deserves to be treated as a
hard architectural constraint, not an optional best practice.

\begin{center}\rule{0.5\linewidth}{0.5pt}\end{center}

\subsection{REFERENCES}\label{references}

{[}1{]} P. Bartlett and M. Wegkamp, ``Classification with a reject
option using a hinge loss,'' \emph{J. Machine Learning Research},
2008.\\
{[}2{]} G. Fumera and F. Roli, ``A theoretical and experimental analysis
of linear combiners for multiple classifier systems,'' \emph{IEEE
TPAMI}, 2005.\\
{[}3{]} B. Zadrozny and C. Elkan, ``Transforming classifier scores into
accurate multiclass probability estimates,'' \emph{KDD}, 2002.\\
{[}4{]} L. Torresani and K.-C. Lee, ``Large margin component analysis,''
\emph{NeurIPS}, 2007.\\
{[}5{]} C. Guo, G. Pleiss, Y. Sun, and K. Q. Weinberger, ``On
calibration of modern neural networks,'' \emph{ICML}, 2017.\\
{[}6{]} Y. Geifman and R. El-Yaniv, ``Selective classification for deep
neural networks,'' \emph{NeurIPS}, 2017.\\
{[}7{]} P. L. Bartlett and M. H. Wegkamp, ``Classification with a reject
option,'' \emph{J. Machine Learning Research}, 2008.\\
{[}8{]} A. Kumar, S. Liang, and T. Ma, ``Verified uncertainty
calibration,'' \emph{NeurIPS}, 2019.\\
{[}9{]} D. Hendrycks and K. Gimpel, ``A baseline for detecting
misclassified and out-of-distribution examples in neural networks,''
\emph{ICLR}, 2017.\\
{[}10{]} K. Lee, K. Lee, H. Lee, and J. Shin, ``A simple unified
framework for detecting out-of-distribution samples and adversarial
attacks,'' \emph{NeurIPS}, 2018.\\
{[}11{]} V. Vovk, A. Gammerman, and G. Shafer, \emph{Algorithmic
Learning in a Random World.} Springer, 2005.\\
{[}12{]} A. N. Angelopoulos and S. Bates, ``A gentle introduction to
conformal prediction and distribution-free uncertainty quantification,''
\emph{arXiv:2107.07511}, 2022.\\
{[}13{]} Y. Geifman and R. El-Yaniv, ``Selectivenet: A deep neural
network with an integrated reject option,'' \emph{ICML}, 2019.\\
{[}14{]} J. Platt, ``Probabilistic outputs for support vector
machines,'' \emph{Advances in Large Margin Classifiers}, 1999.\\
{[}15{]} B. Zadrozny and C. Elkan, ``Obtaining calibrated probability
estimates from decision trees and naive Bayesian classifiers,''
\emph{ICML}, 2001.\\
{[}16{]} S. Liang, Y. Li, and R. Srikant, ``Enhancing the reliability of
OOD image detection in neural networks,'' \emph{ICLR}, 2018.\\
{[}17{]} J. Gama et al., ``A survey on concept drift adaptation,''
\emph{ACM Comput. Surv.}, 2014.\\
{[}18{]} M. Baena-García et al., ``Early drift detection method,''
\emph{ECML-PKDD Workshop}, 2006.\\
{[}19{]} N. Moustafa and J. Slay, ``UNSW-NB15: A comprehensive dataset
for network intrusion detection systems,'' \emph{IEEE MilCIS}, 2015.\\
{[}20{]} M. Tavallaee et al., ``A detailed analysis of the KDD CUP 99
data set,'' \emph{IEEE CISDA}, 2009.

\begin{center}\rule{0.5\linewidth}{0.5pt}\end{center}

\end{document}
